%
% %%%

%%%   Самарский государственный технический университет

%%%   Кафедра прикладной математики и информатики

%%%

%%%   Преамбула для статьи в журнал "Вестник Самарского государственного

%%%   технического университета. Серия: «Физико-математические науки»"

%%%   Ориентирована под MikTEx 2.4 for Win32

%%%

%%%   Хотя сам журнал собирается под GNU/Linux на дистрибутиве TeXLive



\documentclass[11pt,a4paper,twoside]{article}



\usepackage[cp1251]{inputenc}

\usepackage[T2A]{fontenc}

\usepackage[english,russian]{babel}

\usepackage{samgtubib}

\usepackage[dvips]{graphicx}

\usepackage[multiple]{footmisc}



\usepackage{samgtu}

\renewcommand{\VestnikYear}{2009} % Год издания Вестника

\renewcommand{\VestnikNumber}{2\,(19)} % Номер Вестника

\renewcommand{\VestnikMonth}{Ноябрь} % Месяц выхода Вестника

\renewcommand{\VestnikData}{01.11.09} % Дата подписания Вестника в печать

\renewcommand{\VestnikUPL}{26,64} % Усл. печ. л.

\renewcommand{\VestnikUIL}{25,73} % Уч.-изд. л.

\renewcommand{\VestnikRegNumber}{479/09} % Регистрационный номер

\renewcommand{\VestnikOrder}{994} % Номер заказа






%%
% \begin{document}


\renewcommand{\firstpage}{192}
\renewcommand{\lastpage}{196}
%\ShortPartition{Небесная механика и астрономия}{Celestial Mechanics} % Раздел Вестника, в который подаётся статья
% Названия разделов см. на сайте при подаче заявки


%%%%%%%%%%%%%%%%%%%%%%%%%%%%%%%%%%%%%%%%%%%%%%%%%%%%%%%%%%%%%%%%%%%%%%%%%%%%%%%%%%%%
%Информация о статье и об авторах на русском и английском языках

% Номер УДК
\udk{519.245:521.3}
\msc{85-08; 70M20, 65C05}% Коды MSC см. http://www.ams.org/msc/

\title{Сравнительный анализ математических моделей для оценки вероятности столкновения  астероида~Апофис с~Земл\"eй}% Полный заголовок
{Сравнительный анализ математических моделей для оценки вероятности \dots}% Заголовок, который идёт в верхний колонтитул

\author{А.\,Ф.~Заусаев, А.\,Е.~Деревянка}% Автор(ы) для заголовка, если авторов несколько укать \inst
{З\,а\,у\,с\,а\,е\,в~А.\,Ф.,\ Д\,е\,р\,е\,в\,я\,н\,к\,а~А.\,Е.}% Автор(ы) для оглавления и колонтитула через \,

\institute{\samgtu.}


\email{E-mails}{zausaev\_af@mail.ru, andr\_derev@mail.ru} %E-mail's авторов

\otherinf{
\fullauthor{Анатолий Федорович Заусаев}
\regalia{д.ф.-м.н., проф.}\position{профессор}
\dept{каф. прикладной математики и информатики}
\fullauthor{Андрей Евгеньевич Деревянка}
\position{студент}
\dept{каф. прикладной математики и информатики}
}

\keywords{астероид Апофис, вероятность столкновения, элементы орбит, метод Монте-Карло}

\workabstract{Реализованы методы оценки вероятности столкновения астероида Апофис с~Землёй.
Вероятность столкновения оценивается двумя методами: методом Монте-Карло 
и~через отношение промежутка значений элементов орбит, приводящих к~столкновению, к~доверительным интервалам элементов орбит.
Установлено, что основной вклад в~величину оценки вероятности столкновения вносит изменение большой полуоси. 
Получена предполагаемая дата и~оценки вероятности столкновения астероида Апофис с~Землёй.
}

\engtitle{Comparative analysis of mathematical models for~estimating the impact probability of~asteroid~Apophis}

\engauthor{A.\,F.~Zausaev, A.\,E.~Derevyanka}{Z\,a\,u\,s\,a\,e\,v~A.\,F.,
D\,e\,r\,e\,v\,y\,a\,n\,k\,a~A.\,E.}


\enginstitute{\engsamgtu.}% Место работы каждого из авторов на англ. языке
% Если несколько, то так  \enginstitute{ Организация 1 \and Организация 2  \and Организация 3}

\otherenginf{
\fullauthor{Anatoliy F. Zausaev}
\regalia{Dr.\,Sci. (Phys. \& Math.)}
\position{Professor}
\dept{Dept. of Applied Mathematics \& Computer Science}
\fullauthor{Andrey E. Derevyanka}
\position{Student}
\dept{Dept. of Applied Mathematics \& Computer Science}
}

\engabstract{Methods for estimating the impact probability of asteroid Apophis were implemented. The probability of a collision was assessed using two methods: the Monte Carlo method, and the second one, utilizing a ratio between the interval of values of orbital elements, leading to a collision, and the confidence intervals of orbital elements. It was established that the variation in the semimajor axis makes the main contribution to the estimation of the impact probability. The expected date and the estimate of the probability of collision of asteroid Apophis and Earth were computed.
}

\engkeywords{asteroid Apophis, impact probability, orbital elements, Monte Carlo method}

\date{10/IV/2012} % Дата поступления статьи в редакцию.
\revisiondate{24/V/2012} % В окончательном виде

%%%%%%%%%%%%%%%%%%%%%%%%%%%%%%%%%%%%%%%%%%%%%%%%%%%%%%%%%%%%%%%%%%%%%%%%%%%%%%%%%%%%


\maketitle

\small
Астероид Апофис был открыт 19 июня 2004 г. Тогда~же были сделаны первые расчёты его орбиты. 
Тринадцатого апреля 2029~г. в 22:12 по~всемирному времени произойдет сближение с~Землёй, 
которое является наиболее тесным из предсказанных заранее. 
По~данным каталога орбитальной эволюции малых тел Солнечной системы (\mbox{\url{http://smallbodies.ru}}) Апофис  приблизится к~Земле в 2029 году на геоцентрическое расстояние 38,9~тыс.~км. 
Вследствие тесного сближения с~Землёй  орбита астероида существенно изменится, и не~исключена вероятность его столкновения с~Землёй в~будущем. 
Приблизительный диаметр астероида составляет 270~м, поэтому он является потенциально опасным для Земли астероидом.

Вследствие того, что элементы орбит вычисляются с~определёнными погрешностями, важно оценить влияние этих погрешностей на~величину вероятности столкновения этого астероида с~Землёй, а также саму вероятность столкновения. 

Вероятность столкновения астероида с~Землёй в~данной работе оценивалась двумя методами: 
методом Монте-Карло (методом статистических испытаний) и через отношение промежутка значений элементов орбит,  приводящих к~столкновению, к~доверительным интервалам элементов орбит.

Начальные данные элементов орбит на~различные даты наблюдений, использованные в данной работе, взяты из
каталога орбитальной эволюции малых тел Солнечной системы (\mbox{\url{http://smallbodies.ru}}) и~приведены  в~табл.~\ref{derev:tablElemOrb}. Здесь \emph{a} "--- большая полуось, \emph{e} "--- эксцентриситет, \emph{i} "--- наклонение, $\Omega$ --- долгота восходящего узла,
$\omega$ --- аргумент перигелия, \emph{M} "--- средняя аномалия.



Параметры орбиты считаются независимыми нормально распределёнными  случайными величинами. 
В~качестве математического ожидания этих величин использовались данные наблюдений.
Допустимые среднеквадратические отклонения $\sigma$ элементов орбит астероида  [\citen{derev:bib:NASA}] приведены в~табл.~\ref{derev:tablSigmaElem}.

Оценка вероятности столкновения астероида с~Землёй может быть получена  через отношение промежутка значений элементов орбит, приводящих к~столкновению, к~доверительным интервалам элементов орбит.

\begin{table}[!t] \footnotesize
	\center
	\caption{\bf Элементы орбит астероида на различные даты наблюдений}\vspace{-3mm}
	\label{derev:tablElemOrb}

\begin{tabular}{c|c|c|c|c|c|c}
\hline 
\scriptsize
Дата & 
\scriptsize 
\emph{a}, а.е. & 
\scriptsize
\emph{e} & 
\scriptsize
\emph{i}, град. & 
\scriptsize
$\Omega$, град. & 
\scriptsize
$\omega$, град. & 
\scriptsize
\emph{M}, град. \\
\hline
\rule{0mm}{12pt}%
06.03.2006 & 0,92239593 & 0,19104000 & 3,331224 & 204,462302 & 126,355659 & 222,272876 \\
04.01.2010 & 0,92241929 & 0,19121109 & 3,331512 & 204,439306 & 126,424463 & 339,948637 \\
27.08.2011 & 0,92230028 & 0,19107611 & 3,331952 & 204,430424 & 126,424469 & 287,582163 \\
\hline
\end{tabular}

\medskip
\medskip
\setlength{\tabcolsep}{13pt}
\caption{\bf Допустимые среднеквадратические отклонения элементов орбит астероида}\vspace{-3mm}
	\label{derev:tablSigmaElem}
\begin{tabular}[h]{*{5}{c|}c}
			\hline
\scriptsize 
\emph{a}, а.е. & 
\scriptsize
\emph{e} & 
\scriptsize
\emph{i}, град. & 
\scriptsize
$\Omega$, град. & 
\scriptsize
$\omega$, град. & 
\scriptsize
\emph{M}, град. \\
\hline
\rule{0mm}{12pt}%
$ 9,6 \cdot 10^{-9}$&$1,0 \cdot 10^{-7}$&$ 3,0 \cdot 10^{-6}$&$1,5 \cdot 10^{-4}$&$ 1,5 \cdot 10^{-4}$&$ 3,0 \cdot 10^{-6}$\\
			\hline
			\end{tabular}
			\vspace{-3mm}
\end{table}

Для описания движения астероида Апофис использовались дифференциальные уравнения движения  c~учётом гравитационных и~релятивистских эффектов от~Солнца, которые в~гелиоцентрической системе координат имеют следующий вид~[\citen{derev:bib:Brumberg}]:
\begin{multline*}
	\frac{d^{2}X}{dt^{2}}=-k^{2}\left(1+m\right)\frac{X}{r^{3}}+\sum_{i}k^{2}m_{i}\Bigl(\frac{X_{i}-X}{\Delta^{3}_{i}}-\frac{X_{i}}{r_{i}^{3}}\Bigr)+\\
	+\frac{k^{2}}{c^{2}}\Bigl(\left(4-2\alpha\right)\frac{k^{2}}{r^{2}}X-\left(1+\alpha\right)\frac{\dot r^{2}}{r^{3}}X+3\alpha\frac{(X\dot X)^{2}}{r^{5}}X+\left(4-2\alpha\right)\frac{(X\dot X)}{r^{3}}\dot X\Bigr).
\end{multline*}
Здесь $X$ "--- матрица-столбец с~элементами $x$, $y$, $z$; 
$X_{i}$ "--- матрица--столбец с~элементами $x_{i}$, $y_{i}$, $z_{i}$;
$m$, $x$, $y$, $z$ "--- масса и гелиоцентрические координаты астероида;
 $m_{i}$, $x_{i}$, $y_{i}$, $z_{i}$ "--- массы и гелиоцентрические координаты больших планет;
 $r^{2}=x^{2}+y^{2}+z^{2}$,
 $\Delta_i=\left(x_{i}-x\right)^{2}+\left(y_{i}-y\right)^{2}+\left(z_{i}-z\right)^{2}$,
 $r^{2}_{i}=x^{2}_{i}+y^{2}_{i}+z^{2}_{i}$;
 $\dot X$ "--- матрица-столбец с~элементами $\dot x$, $\dot y$, $\dot z$;
 $k$ "--- постоянная Гаусса; $c$ "--- скорость света, $\alpha$ "--- параметр, характеризующий выбор системы координат.
 Случай $\alpha=1$ соответствует стандартным координатам, $\alpha =0$ "--- гармоническим координатам. 


Численное интегрирование уравнений движения астероида методом Эверхарта 27~порядка  [\citen{derev:bib:Everhart,derev:bib:Zausaevs}]
для различных вариантов начальных данных в~пределах вышеуказанных отклонений проведено на~интервале времени с~2006 по~2040~гг.

Начальные данные элементов орбит астероида на~различные моменты оскулляции варьировались в~пределах их доверительных интервалов. При~варьировании каждого элемента  оценивалось минимальное сближение и~дата сближения астероида с~Землёй. 
Считалось, что происходит столкновение астероида с~Землёй,  если расстояние между центрами Земли и астероида меньше радиуса Земли.

На~основании проведённых расчётов установлено, что столкновение между астероидом Апофис и~Землёй может произойти только при изменении большой полуоси  в~пределах доверительного интервала ($\emph{a}-3\sigma_{\emph{a}}; \emph{a}+3\sigma_{\emph{a}}$).
Здесь за $\emph{a}_{1}$ и $\emph{a}_{2}$ приняты минимальное и максимальное значения большой полуоси астероида, при~которых возможны его столкновения с~Землёй, а $\sigma_{\emph{a}}$ "--- допустимое отклонение для большой полуоси астероида.
В~результате расчётов определена предполагаемая дата столкновения Апофиса с~Землёй --- 13.04.2037.

Вероятность столкновения астероида Апофис может быть оценена через отношение  величины $\left|a_{1}-a_{2}\right|$ к~величине доверительного интервала:
$$
	P=|a_{1}-a_{2}|/(6\sigma_{a}).
$$
По этой формуле получены вероятности столкновения астероида Апофис с~Землёй  на дату 13.04.2037 на основании различных дат наблюдений, которые приведены  в~табл.~\ref{derev:tablRes}.

При~совместном изменении значений большой полуоси \emph{a} и эксцентриситета \emph{e}  установлено, что эксцентриситет оказывает дополнительное влияние 
на~величину сближения с~Землёй, но не является определяющим, так как изменение только эксцентриситета
 в~пределах доверительно интервала ($\emph{e}-3\sigma_{\emph{e}};\emph{e}+3\sigma_{\emph{e}}$) 
не приводит к~столкновению.

Суть метода Монте-Карло применительно к~рассматриваемой задаче заключается в~том, что начальные условия движения астероида (виртуальные астероиды) выбираются  в~пределах области возможных начальных значений, определённой на~основе имеющегося набора наблюдений. 
Под~этой областью подразумевается некоторое множество  из~шестимерного фазового пространства орбитальных элементов или 
кинематических параметров "--- декартовых координат радиус-вектора и компонентов вектора скорости. 
Далее движение виртуальных астероидов  с~заданной точностью численного интегрирования прослеживается на~некотором промежутке 
времени до~тех~пор, пока они не столкнутся с~планетой или не пролетят мимо неё. 
Формально считается, что столкновение с~планетой происходит в~том случае, если при сближении расстояние между центрами планеты и~астероида становится меньше  радиуса планеты. 
При большом числе испытаний отношение числа начальных условий, приводящих к~столкновению, к~общему числу испытаний может рассматриваться как  вероятность столкновения [\citen{derev:bib:Jeleznov}]. 
Сходимость метода Монте-Карло является сходимостью по вероятности [\citen{derev:bib:Ermakov}]. 
Метод Монте-Карло  прост, но чрезвычайно трудоёмок: чем меньше вероятность  события, тем больше начальных условий движения необходимо использовать для  получения достоверной оценки вероятности.

На~различные моменты оскуляции производилось интегрирование уравнений движения 30\,000 виртуальных астероидов с~начальными данными, являющимися независимыми нормально распределенными случайными величинами. 
Уравнения движения этих астероидов интегрировались методом Эверхарта 27~порядка 
на~интервале времени с~2006 по~2040 гг. 
Соударение засчитывалось, если величина сближения была меньше радиуса Земли (6378~км).

Вероятности столкновения астероида Апофис с Землёй на дату 13.04.2037, полученные методом Монте-Карло также представлены в~табл.~\ref{derev:tablRes}. Видно, что значения вероятностей, полученные различными методами на~различные моменты оскуляции,  имеют одинаковый порядок и~расположены достаточно близко друг к~другу.

\begin{table}[t!]
	\center \footnotesize
	\caption{\textbf{Вероятности столкновения астероида Апофис с~Землёй на дату 13.04.2037   на основании различных дат наблюдений }}
\vspace{-3mm}
	\label{derev:tablRes}
	\setlength{\tabcolsep}{16pt}
\begin{tabular}[h]{c|c|c}\hline
\scriptsize
Дата &
\scriptsize
Метод Монте-Карло& 
\scriptsize
Отношение опасного промежутка к $6\sigma_{a}$\\\hline
\rule{0mm}{12pt}%
			06.03.2006&$5,33 \cdot 10^{-4}$& {$3,41 \cdot 10^{-4}$}\\
			04.01.2010&$5,67 \cdot 10^{-4}$& {$3,74 \cdot 10^{-4}$}\\
			27.08.2011&$3,67 \cdot 10^{-4}$& {$3,99 \cdot 10^{-4}$}\\\hline
			\end{tabular}
\vspace{-5mm}
\end{table}

Полученные значения вероятности столкновения астероида Апофис с~Землёй на дату 13.04.2037 не противоречат результатам других исследователей. Так, вероятность столкновения, предложенная в~[\citen{derev:bib:NASA}] равна $4,5\cdot 10^{-4}$, а вероятность, предложенная в~[\citen{derev:bib:Smirnov}], равна $2,5\cdot 10^{-5}$.
В~работах [\citen{derev:bib:NASA}, \citen{derev:bib:Smirnov}] использовался метод Монте-Карло. 


\thanks{Работа выполнена в рамках государственного задания высшим учебным заведениям в~части проведения
научно"=исследовательских работ (проект 2.534.2011).}
%% Благодарности, гранты и прочее

\begin{thebibliography}{9}

\Bibitem{derev:bib:NASA}
\by Giorgini~J.\,D., Benner~L.\,A.\,M., Ostro~S.\,J., Nolan~M.\,C., Busch~M.\,W.
\paper Predicting the Earth encounters of (99942) Apophis
\jour Icarus
\vol 193
\issue 1
\yr 2008
\pages 1--19


\RBibitem{derev:bib:Brumberg}
\book Релятивистская небесная механика
\by Брумберг~В.\,А.
\publaddr М.
\publ Наука
\yr 1972
\totalpages 382
\misctransl
\book Relativistic celestial mechanics
\by Brumberg~V.\,A.
\publaddr Moscow
\publ Nauka
\yr 1972
\totalpages 382

\Bibitem{derev:bib:Everhart}
\paper Implicit Single-Sequence Methods for Integrating Orbits
\by Everhart~E.
\jour Cel. Mech.
\yr 1974
\vol 10
\issue 1
\pages 35--55

\RBibitem{derev:bib:Zausaevs} 
\paper Применение модифицированного метода Эверхарта  для решения задач небесной механики
\by Заусаев~А.\,Ф., Заусаев~А.\,А.
\jour Матем. моделирование
\yr 2008
\publaddr М.
\vol 20
\issue 11
\pages 109--114
\misctransl
\by Zausaev~A.\,F., Zausaev~A.\,A.
\paper Employment of the modification Everhart's method for solution of problems of celestial mechanics
\jour Matem. Mod.
\yr 2008
\vol 20
\issue 11
\pages 109--114
		
\RBibitem{derev:bib:Jeleznov}
\paper Влияние корреляционных связей  между оцениваемыми по~наблюдениям орбитальными параметрами астероида на~результаты определения вероятности его столкновения с~планетой методом Монте-Карло
\by Железнов~Н.\,Б.
\jour Астрон. вестн.
\yr 2010
\vol 44
\issue 2
\pages 150--157
\transl англ. пер.:~
\paper The influence of~correlations  between asteroid orbital parameters estimated by~observations on~the~probability of~asteroid encounter with a~planet calculated by~Monte~Carlo method
\by Zheleznov~N.\,B.
\jour Solar System Research
\yr 2010
\vol 44
\issue 2
\pages 150--157


\RBibitem{derev:bib:Ermakov}
\book Метод Монте-Карло и смежные вопросы
\by Ермаков~С.\,М.
\yr 1975
\publaddr М.
\publ Наука
\totalpages 472
\misctransl
\book The Monte Carlo method and related problems
\by Ermakov~S.\,M.
\yr 1975
\publaddr Moscow
\publ Nauka
\totalpages 472

\RBibitem{derev:bib:Smirnov}
\by Смирнов~Е.\,А.
\inbook  Астрономия и всемирное наследие через время и континенты
\bookinfo Тр. международн. конф. \nofrills
\paper Использование интервальной арифметики при прогнозировании орбит малых тел
\publaddr Казань
\publ Казан. гос. ун-т
\pages 101--102
\yr 2009
\misctransl
\by Smirnov E.\,A.
\inbook Astronomy and World Heritage across Time and Continents
\bookinfo Proc. of Internat. Conf. \nofrills
\paper The use of~interval arithmetic in~predicting the~orbits of~small bodies
\publaddr Kazan
\publ Kazan. Gos. Un-t
\pages 101--102
\yr 2009


\end{thebibliography}

\noindent
\hskip6.5cm \thedate \par\noindent
\hskip6.5cm\therevdate

\vspace{-3mm}

%\newpage
%
%\chead[
%	\fancyplain{}{
%	\scriptsize \sl Z\,a\,u\,s\,a\,e\,v~A.\,F.,
%	D\,e\,r\,e\,v\,y\,a\,n\,k\,a~A.\,E.
%	}]
%	{
%	\fancyplain{}
%	{\scriptsize \sl  Comparative analysis of mathematical models for estimating \dots}
%	}
%
%	
\makeengtitle

 \newpage
%
% \tableofengcontents
%%
%\end{document}